% Industry research CV, adapted from Star Rover by Subidit (CC BY 4.0).
% See TEMPLATE-LICENSE.md.
\documentclass[10pt]{article}

\renewcommand{\familydefault}{\sfdefault}
\usepackage[a4paper,margin=0.78in]{geometry}
\usepackage{xcolor}
\definecolor{accent}{HTML}{141E61}
\setcounter{secnumdepth}{0}
\pdfgentounicode=1

\usepackage{enumitem}
\setlist[itemize]{nosep,left=0pt .. 1.5em}
\setlist[enumerate]{left=0pt..1.5em,itemsep=2pt}
\setlist[description]{itemsep=0pt}

\usepackage{xhfill}
\usepackage{soul}
\newcommand{\midrule}[1]{%
  \leavevmode\xrfill[.5ex]{1pt}[accent]~\so{#1}~\xrfill[.5ex]{1pt}[accent]}
\usepackage{titlesec}
\titlespacing{\section}{0pt}{*3}{*0.8}
\titlespacing{\subsection}{0pt}{*2}{*0}
\titlespacing{\subsubsection}{0pt}{*0}{*0}
\titleformat{\section}{\large\bfseries\uppercase}{}{}{\midrule}
\titleformat*{\subsection}{\large\bfseries\scshape}
\newcommand{\aux}[1]{$|$ \space {\normalfont\itshape #1}}
\newcommand{\rside}[1]{\hfill {\normalfont\color{accent} #1}}

\usepackage{lastpage}
\usepackage{fancyhdr}
\pagestyle{fancy}
\fancyhf{}
\renewcommand{\headrulewidth}{0pt}
\fancyfoot[C]{\small\color{gray}Pingshi Yu -- Industry Research CV -- Page \thepage\ of \pageref*{LastPage}}
\usepackage[bookmarks=false,hidelinks]{hyperref}
\newcommand{\linkicon}{{\color{gray}\footnotesize[link]}}

\begin{document}

\begin{center}
  {\fontsize{30}{34}\fontseries{heavy}\selectfont\color{accent}
    PINGSHI (JACOB) YU}\\[4pt]
  {\large Research Scientist / Research Engineer --- Programming Languages and Software Reliability}\\[5pt]
  London, UK \quad
  \href{mailto:p.yu22@imperial.ac.uk}{p.yu22@imperial.ac.uk} \quad
  \href{https://pingshiyu.github.io}{pingshiyu.github.io}\\[2pt]
  \href{https://github.com/pingshiyu}{GitHub: pingshiyu} \quad
  \href{https://www.linkedin.com/in/jacob-pingshi-yu/}{LinkedIn} \quad
  \href{https://scholar.google.com/citations?user=hbgX7rEAAAAJ}{Google Scholar}
\end{center}

\section{Profile}

Final-year Computer Science PhD researcher combining programming-languages theory with practical systems implementation. In my research, I build methods and tools for finding compiler bugs and reasoning about software correctness. This work spans fuzzing, composable compiler semantics, property-based testing, complexity theory, and automated verification. Prior industry experience includes programming-language research, quantitative engineering for an FX market-making platform, quantum software, and data engineering.

\section{Education}

\subsection{Imperial College London \aux{PhD in Computer Science} \rside{2022--present}}
Full doctoral scholarship. Research on compiler testing, composable semantics, property-based testing, and the complexity of randomised testing. Supervised by Prof. Alastair F. Donaldson and Prof. Nicolas Wu.

\subsection{University of Oxford, St Anne's College \aux{MMathCompSci} \rside{2016--2020}}
Mathematics and Computer Science; first-class master's year (80.3\%). Dissertation: \textit{Systematic Pruning Methods for the Enumeration of Realistic Molecules}.

\section{Research Highlights}

\subsection{Compiler Testing and Composable Semantics \rside{Imperial College London}}
\begin{itemize}
  \item Developed \textit{Ratte}, a semantics-driven approach to fuzzing multi-level compilers for miscompilations; published and presented at ASPLOS 2025.
  \item Researched effect-handler-based semantics for composable compiler IRs, connecting executable language semantics with automated compiler validation.
  \item Contributed to \textit{RustSmith}, a random differential testing tool for the Rust compiler, published at ISSTA 2023.
\end{itemize}

\subsection{Theory of Randomised Testing \rside{Imperial College London}}
\begin{itemize}
  \item Developed a complexity-theoretic account of randomised testing, linking practical testing strategies with formal questions about bug-finding difficulty; released on arXiv in 2026.
\end{itemize}

\subsection{Effect Handlers for Cangjie \rside{Huawei Research}}
\begin{itemize}
  \item Researched novel applications of effect handlers---a programming-language feature---for software testing in Huawei's open-source Cangjie programming language.
\end{itemize}

\section{Industry and Research Experience}

\subsection{Huawei Research \rside{Edinburgh, UK}}
\subsubsection{Research Intern, Programming Languages Lab \rside{2025--2026}}
\begin{itemize}
  \item Researched novel applications of effect handlers (a programming-language feature) for software testing in Huawei's open-source Cangjie programming language.
\end{itemize}

\subsection{Credit Suisse \rside{London, UK}}
\subsubsection{Software Engineer / Quantitative Strategist \rside{2020--2022}}
\begin{itemize}
  \item Developed components of a production foreign-exchange market-making platform.
  \item Led delivery of NLP-derived market-data pipelines into the pricing engine, coordinating technical direction across quantitative, support, and data teams.
  \item Built services for quantitative research, model validation, release management, performance monitoring, and developer productivity.
  \item Worked with Python, C++, Docker, Jenkins, JavaScript, Boost, SQL, and Elasticsearch.
\end{itemize}

\subsection{Cambridge Quantum Computing (now Quantinuum) \rside{Cambridge, UK}}
\subsubsection{Research Intern \rside{2020}}
\begin{itemize}
  \item Surveyed ZX-calculus research and implemented a scalable graphical calculus in C++ for representing, reasoning about, and simplifying large quantum circuits.
\end{itemize}

\subsection{Earlier Engineering and Data Experience \rside{2017--2019}}
\begin{itemize}
  \item \textbf{Credit Suisse:} delivered two production projects and identified an optimisation yielding approximately $900\times$ faster bulk SQL inserts in an internal Perl API.
  \item \textbf{Arm:} initiated and delivered the automated data-processing backend for a company-wide capacity-planning project.
  \item \textbf{University of Oxford DPIR:} built a Python pipeline to process roughly 10,000 statutory instruments and designed validation materials for sentiment analysis.
\end{itemize}

\section{Publications}

\begin{enumerate}
  \item \textbf{Pingshi Yu}, Chengsong Tan, Nicolas Wu, and Alastair F. Donaldson. ``Complexity Theory of Randomised Testing.'' arXiv, 2026. \href{https://arxiv.org/abs/2607.11811}{Paper \linkicon}
  \item \textbf{Pingshi Yu}, Nicolas Wu, and Alastair F. Donaldson. ``Ratte: Fuzzing for Miscompilations in Multi-Level Compilers Using Composable Semantics.'' \textit{ASPLOS}, 2025. \href{https://doi.org/10.1145/3676641.3716270}{DOI \linkicon}
  \item Mayank Sharma, \textbf{Pingshi Yu}, and Alastair F. Donaldson. ``RustSmith: Random Differential Compiler Testing for Rust.'' \textit{ISSTA}, 2023. \href{https://doi.org/10.1145/3597926.3604919}{DOI \linkicon}
  \item \textbf{Pingshi Yu}. ``Reasoning about MLIR Semantics through Effects and Handlers.'' \textit{ECOOP/ISSTA Doctoral Symposium}, 2023. \href{https://doi.org/10.1145/3597926.3605239}{DOI \linkicon}
  \item \textbf{Pingshi Yu}, Alistair J. Sterling, and Jotun Hein. ``A Novel Automated Screening Method for Combinatorially Generated Small Molecules.'' \textit{Journal of Chemical Information and Modeling}, 2021. \href{https://doi.org/10.1021/acs.jcim.0c01462}{DOI \linkicon}
\end{enumerate}

\section{Technical Skills}

\begin{description}[style=multiline,leftmargin=9.5em]
  \item[Programming] Haskell, Python, C++; experience with Agda, Java, Scala, Perl, and MATLAB.
  \item[Research areas] Compilers, fuzzing, property-based testing, programming-language semantics, effect systems and handlers, algorithms, automated verification, complexity theory.
  \item[Engineering] Git, Unix, Docker, Jenkins, SQL, Boost, Pandas, TensorFlow, Elasticsearch, \LaTeX.
  \item[Languages] Chinese (native), English (native), German (basic).
\end{description}

\section{Recognition}

\begin{itemize}
  \item Full Doctoral Scholarship, Imperial College London, 2022--2027.
  \item Poster Presentation Award, Imperial Computing Conference, 2024.
  \item Invited research presentation at ETH Z\"urich, 2024; paper presentations at ASPLOS 2025 and ISSTA 2023.
  \item Exhibition Awards, St Anne's College, University of Oxford, 2019 and 2020.
\end{itemize}

\vspace{4pt}
\begin{center}{\small\color{gray}Last updated August 2026}\end{center}

\end{document}
